% CPSY 1291 -- Recitation handout 3
% Compile:  latexmk -lualatex handout-03-matrices-fitting-knobs.tex
\documentclass[11pt]{article}
\usepackage{cpsy1291-handout}

\begin{document}

\handouthead{3}{Matrices, curve fitting, and methods with knobs}
  {The week before Assignment 1 is due}
  {Week 3 \ $\cdot$\ hold this after Lecture 4, before Tue 9/29}

\section*{What this session is for}
\addcontentsline{toc}{section}{What this session is for}

Recitations 1 and 2 gave you arrays, shapes, distances, dissimilarity matrices
and PCA. Assignment 1 needs four more things, and none of them is taught in
lecture:

\begin{enumerate}
  \item \textbf{matrix multiplication} as an operation you can predict the shape
        and the meaning of, rather than a symbol you try until it runs;
  \item \textbf{fitting a curve to data} and comparing two candidate curves
        honestly (Part 2 asks you to test Shepard's law);
  \item \textbf{rank correlation}, and why comparing two dissimilarity matrices
        uses it rather than Pearson (Part 2k);
  \item \textbf{methods with knobs} --- $t$-SNE's perplexity, every method's
        \texttt{random\_state} --- and how to report a result that depends on
        them (Part 4).
\end{enumerate}

\begin{keybox}
\ask{The rule this week.} A method with a knob does not give you a result. It
gives you a \emph{family} of results, one per knob setting. Your job is to say
which features of that family survive the knob and which do not --- and only the
survivors are findings.
\end{keybox}

\section{Matrix multiplication, and how to predict it}

\subsection{The shape rule, and why it is the whole thing}

If $\boldsymbol{A}$ is $(m \times k)$ and $\boldsymbol{B}$ is $(k \times n)$,
then $\boldsymbol{A}\boldsymbol{B}$ is $(m \times n)$. The inner dimensions must
match and they \emph{vanish}; the outer dimensions survive.

\begin{lstlisting}
A = np.zeros((120, 4096))     # 120 images x 4096 units
B = np.zeros((4096, 10))      # 4096 units  x 10 categories
(A @ B).shape                 # (120, 10)   -- 4096 cancelled
\end{lstlisting}

Say the cancelled dimension out loud: ``\emph{summed over units}''. That is what
a matrix product is --- a sum over the shared index --- and naming it tells you
immediately whether the operation means anything.

\begin{keybox}
\ask{\texttt{@} is not \texttt{*}.} \texttt{A * B} multiplies element by element
and broadcasts. \texttt{A @ B} contracts the shared index. They agree on
nothing, and NumPy will happily run the wrong one when the shapes permit. If you
wrote \texttt{*} and got an answer with the shape you expected, check that you
did not mean \texttt{@}.
\end{keybox}

\subsection{The three products you will actually write}

Let $\boldsymbol{X}$ be $(n \times p)$: $n$ stimuli down the rows, $p$ features
across the columns. Then

\begin{center}
\begin{tabular}{@{}lll@{}}
\toprule
Expression & Shape & What it holds \\
\midrule
\texttt{X @ X.T} & $(n \times n)$ & every pair of \emph{stimuli}, dotted \\
\texttt{X.T @ X} & $(p \times p)$ & every pair of \emph{features} $\rightarrow$ covariance \\
\texttt{X @ v}   & $(n,)$         & every stimulus projected onto direction $\boldsymbol{v}$ \\
\bottomrule
\end{tabular}
\end{center}

The first is where dissimilarity matrices come from. The second, after
centering and dividing by $n-1$, \emph{is} the covariance matrix whose
eigenvectors PCA returns. The third is what \texttt{pca.transform} does, once
per component. Three lines, and they are most of the assignment.

\subsection{Transpose, and the mistake it causes}

\texttt{X.T} flips rows and columns. The failure mode is not that it errors ---
it is that on a square matrix it does not, and you silently compute the
covariance of your stimuli instead of your features. Whenever you write
\texttt{.T}, write the resulting shape in a comment. It costs three seconds and
it is the single highest-yield habit in this course.

\begin{keybox}
\ask{Order matters.} $\boldsymbol{A}\boldsymbol{B} \neq
\boldsymbol{B}\boldsymbol{A}$, even when both are defined. Matrix
multiplication composes operations, and composing in the other order is a
different operation --- rotate-then-stretch is not stretch-then-rotate.
\end{keybox}

\section{Reading an eigenvalue spectrum}

Recitation 2 said what an eigenvector is: a direction the covariance matrix
stretches without rotating, and PC1 is the most-stretched one. Here is the part
you need for Part 3.

\texttt{pca.explained\_variance\_ratio\_} is not a number, it is a
\textbf{curve}: eigenvalue 1, eigenvalue 2, \dots each divided by their total.
Sorted, always decreasing. What the curve looks like is the finding.

\begin{lstlisting}
ratio = pca.explained_variance_ratio_
plt.plot(np.cumsum(ratio))          # how much you have kept by component k
plt.axhline(0.9, ls=':')            # a threshold you must choose and defend
\end{lstlisting}

\begin{itemize}
  \item A spectrum that \textbf{falls off a cliff} means a genuinely
        low-dimensional representation: a handful of directions carry it.
  \item A spectrum that \textbf{decays slowly} means variance is spread across
        many directions, and no small number of components summarizes it.
  \item ``$k$ components explain $90\%$'' is meaningless without saying
        \emph{out of how many}, and \emph{on how many items}.
\end{itemize}

That last point is the trap. If your stimuli are $51$ objects photographed at
$24$ angles each, rotating one object traces a smooth curve through feature
space: you have added $24\times$ the points but hardly any new directions. The
representation looks compact, and the number changes if you change the number of
viewpoints while the network stays byte-for-byte identical.

\begin{keybox}
\ask{Dimensionality is a property of the data \emph{and} the stimulus set.} Never
report it as a property of the network alone.
\end{keybox}

\section{Fitting a curve, and comparing two}

\subsection{The two calls}

For a polynomial, NumPy is enough:

\begin{lstlisting}
c    = np.polyfit(x, y, deg=1)      # coefficients, highest power first
yhat = np.polyval(c, x)
\end{lstlisting}

For anything else --- an exponential, a Gaussian --- you write the function and
let SciPy find the parameters:

\begin{lstlisting}
from scipy.optimize import curve_fit

def expo(d, a, b):  return a * np.exp(-b * d)

p, _ = curve_fit(expo, d, s, p0=[1.0, 1.0])   # p0 = a starting guess
yhat = expo(d, *p)
\end{lstlisting}

\texttt{curve\_fit} minimizes squared error by local search, so it needs a
starting guess and can land in a bad local minimum. If a fit looks absurd, try a
different \texttt{p0} before you conclude anything about the data.

\subsection{Comparing candidate curves}

\[
R^2 \;=\; 1 - \frac{\sum_i (y_i - \hat{y}_i)^2}{\sum_i (y_i - \bar{y})^2}
\]

\begin{lstlisting}
def r2(y, yhat):
    return 1 - ((y - yhat)**2).sum() / ((y - y.mean())**2).sum()
\end{lstlisting}

Three cautions, all of which bite in Part 2:

\begin{enumerate}
  \item \textbf{$R^2$ is a ratio, so the denominator matters.} Bin your data
        before fitting and $R^2$ rises with the fitted curve unchanged --- you
        removed noise from the denominator, not error from the model. Two $R^2$
        values are comparable only against the same target.
  \item \textbf{More parameters never fit worse.} A curve with three free
        parameters will beat one with two on the same data, always. That is
        arithmetic, not evidence. Compare models with the \emph{same} number of
        parameters, or say plainly that you did not.
  \item \textbf{An $R^2$ can be negative}, if your fit is worse than predicting
        the mean. This is a real outcome, not a bug.
\end{enumerate}

\begin{keybox}
\ask{What Shepard's law claims.} That perceived similarity falls
\emph{exponentially} with distance in psychological space --- $s = a\,e^{-bd}$,
not $s = a - bd$ and not a Gaussian. Testing it means fitting all three and
comparing, which is exactly the machinery above.
\end{keybox}

\section{Pearson, Spearman, and the upper triangle}

\subsection{Which correlation, and why}

\textbf{Pearson} measures how close the relationship is to a straight line.
\textbf{Spearman} replaces every value by its rank and then computes Pearson on
the ranks, so it measures whether the relationship is \emph{monotonic} --- when
one goes up, does the other?

\begin{lstlisting}
from scipy.stats import pearsonr, spearmanr
r, p   = pearsonr(a, b)
rho, p = spearmanr(a, b)
\end{lstlisting}

Comparing two dissimilarity matrices uses Spearman, and the reason is
substantive. A network's distances and a human's similarity judgments are not
on the same scale and there is no reason to expect a linear relation between
them. What you want to ask is: \emph{do the two systems order the pairs the same
way?} That question is Spearman's, and Pearson would answer a different one and
report a lower number for a reason that has nothing to do with the models.

\subsection{Never correlate the whole matrix}

A dissimilarity matrix is symmetric with a zero diagonal, so more than half its
entries are duplicates or structural zeros. Feeding all $n^2$ of them into a
correlation inflates the sample and counts every pair twice. Take the upper
triangle:

\begin{lstlisting}
iu = np.triu_indices(D1.shape[0], k=1)   # k=1 skips the diagonal
rho, p = spearmanr(D1[iu], D2[iu])
\end{lstlisting}

\begin{keybox}
\ask{The $p$-value here is wrong, and you should say so.} The $n(n-1)/2$ pairs
are not independent --- they are built from only $n$ stimuli, so moving one
stimulus moves a whole row and column at once. The correlation is meaningful;
its $p$-value is not, and the honest move is to report the correlation and note
the dependence. (The standard fix is a permutation test over \emph{stimuli},
which you will meet in Recitation 8.)
\end{keybox}

\section{Methods with knobs}

\subsection{\texttt{random\_state}: reproducible is not the same as meaningful}

MDS and $t$-SNE both start from a random initialization and optimize. Two runs
give two different pictures.

\begin{lstlisting}
Y = TSNE(n_components=2, perplexity=30, random_state=0).fit_transform(X)
\end{lstlisting}

Setting \texttt{random\_state} makes the run \textbf{reproducible}: you will get
the same picture tomorrow. It does not make the picture \textbf{meaningful}: you
have picked one member of a family and hidden the rest. The way to tell the
difference is to run several seeds and look at what stays put.

\begin{keybox}
\ask{The test to apply.} Rerun with three different seeds. Structure that
appears in all three is a candidate finding. Structure that appears in one is a
property of that seed, and reporting it is an error --- a reproducible one.
\end{keybox}

\subsection{Perplexity, and what a $t$-SNE map is not}

$t$-SNE's \texttt{perplexity} sets roughly how many neighbors each point tries to
stay near --- small values preserve local detail, large values preserve coarser
structure. Sweep it and the picture changes substantially.

Three things a $t$-SNE map does \emph{not} tell you, all of which students read
off it anyway:

\begin{itemize}
  \item \textbf{Distances between clusters are not distances.} Two clusters
        drawn far apart are not more different than two drawn close together.
  \item \textbf{Cluster sizes are not sizes.} The algorithm expands sparse
        regions and compresses dense ones by design.
  \item \textbf{Apparent clusters are not evidence of clusters.} $t$-SNE will
        produce well-separated blobs from data with no cluster structure at all,
        at some perplexities.
\end{itemize}

What it \emph{is} good for: seeing whether points that you already have labels
for land together. That is a legitimate and useful question, and it is
essentially the only one the map answers reliably.

\subsection{What two dimensions cost}

Any method that puts high-dimensional data on a page must throw something away.
The question worth asking --- and the one Part 4 asks --- is what. A usable
measure is \textbf{neighborhood preservation}: for each point, what fraction of
its $k$ nearest neighbors in the original space are still among its $k$ nearest
neighbors in the embedding?

\begin{lstlisting}
def neighbor_overlap(D_hi, D_lo, k=10):
    n   = D_hi.shape[0]
    hi  = np.argsort(D_hi, axis=1)[:, 1:k+1]     # skip self at position 0
    lo  = np.argsort(D_lo, axis=1)[:, 1:k+1]
    return np.array([len(set(hi[i]) & set(lo[i])) / k for i in range(n)])
\end{lstlisting}

This returns one number per point, so you can also \emph{map} it --- color each
point by how well its neighborhood survived --- and see which regions of the
embedding you are allowed to trust.

\section{A recipe for when it breaks}

\begin{enumerate}
  \item \textbf{Read the last line of the traceback first.} It names the error;
        everything above it is the route taken to get there.
  \item \textbf{Then find the last line that is \emph{your} file.} The frames
        below it are inside a library and are almost never the bug.
  \item \textbf{Print shapes, not values.} \texttt{print(X.shape, Y.shape)}
        resolves most of them outright.
  \item \textbf{Check for \texttt{nan} early.} \texttt{np.isnan(X).sum()} ---
        one \texttt{nan} propagates through a whole matrix silently. Correlation
        distance produces them for any constant row (a dead unit), which is why
        Part 2c makes you look.
  \item \textbf{Shrink the input.} Run on the first 5 rows. If the bug survives,
        you can now inspect every number by hand; if it does not, it is about
        scale, not logic.
\end{enumerate}

\section{Exercises}

\ask{1.} \texttt{X} is $(120, 4096)$ and \texttt{W} is $(4096, 64)$. What are
the shapes of \texttt{X @ W}, \texttt{W.T @ X.T}, and \texttt{X @ X.T}? Which
one is a dissimilarity-matrix-shaped object?

\ask{2.} You compute $R^2 = 0.84$ for an exponential fit. A paper reports
$0.96$ for the same law on the same kind of data. Name two things that could
differ between you and the paper that would produce this gap \emph{without}
either model being better.

\ask{3.} You compare a network RDM with a human RDM and get Spearman
$\rho = 0.42$, $p < 10^{-9}$ over $7{,}140$ pairs. What is wrong with the
$p$-value, and what would you report instead?

\ask{4.} A $t$-SNE map shows two clean, widely separated clusters. Give two
different explanations, one of which is a real finding and one of which is not,
and state the single check that separates them.

\ask{5.} A colleague reports that a network's representation is
``$8$-dimensional, since $8$ components explain $90\%$ of the variance.'' Their
stimulus set is $40$ objects at $30$ viewpoints. What is the objection?

\subsection*{Answers}

\ask{1.} $(120, 64)$; $(64, 120)$; $(120, 120)$. The last one --- it is square
and indexed by stimuli on both sides.

\ask{2.} (i) They may have \emph{binned} the data before fitting, which shrinks
the denominator of $R^2$ without improving the fit. (ii) They may have fitted a
model with more free parameters. (Also acceptable: a different stimulus set with
more spread in $y$, or averaging over subjects before fitting.)

\ask{3.} The $7{,}140$ pairs come from only $120$ stimuli and are therefore not
independent, so the $p$-value's assumption fails and the number is far too
small. Report $\rho$ with the number of \emph{stimuli}, and get a $p$-value from
a permutation test that shuffles stimulus labels.

\ask{4.} Real: the two groups genuinely occupy separate regions of the
high-dimensional space. Not real: $t$-SNE manufactured the separation at this
perplexity, or this seed. The check: rerun across several seeds \emph{and}
several perplexities, and see whether the split survives.

\ask{5.} $1{,}200$ images, but only $40$ distinct objects; the $30$ viewpoints of
one object trace a smooth low-dimensional curve rather than adding independent
directions. The $90\%$ figure is partly a fact about how the stimulus set was
built, and it would move if they photographed $10$ viewpoints instead --- with
the network unchanged.

\section*{Cheat sheet}
\addcontentsline{toc}{section}{Cheat sheet}

\begin{center}
\begin{tabular}{@{}p{0.44\textwidth}p{0.5\textwidth}@{}}
\toprule
\textbf{You want} & \textbf{You write} \\
\midrule
every pair of stimuli, dotted        & \texttt{X @ X.T} \\
covariance of the features           & \texttt{np.cov(X, rowvar=False)} \\
project onto a direction             & \texttt{X @ v} \\
how much each component explains     & \texttt{pca.explained\_variance\_ratio\_} \\
fit a line                           & \texttt{np.polyfit(x, y, 1)} \\
fit your own function                & \texttt{curve\_fit(f, x, y, p0=...)} \\
goodness of fit                      & \texttt{1 - sse / sst} \\
monotonic agreement                  & \texttt{spearmanr(a, b)} \\
upper triangle of an RDM             & \texttt{np.triu\_indices(n, k=1)} \\
same picture twice                   & \texttt{random\_state=0} \\
count the \texttt{nan}s              & \texttt{np.isnan(X).sum()} \\
$k$ nearest neighbors of each row    & \texttt{np.argsort(D, axis=1)[:, 1:k+1]} \\
\bottomrule
\end{tabular}
\end{center}

\end{document}
